\section[Conclusão]{Conclusão}

Sendo minha terceira experiência na Ages, o projeto VanGO foi responsável por uma mudança
importante de responsabilidades, quando comparado com as etapas anteriores. Além de
atuar tecnicamente no \textit{back-end}, também precisei acompanhar decisões que diziam respeito a
arquitetura, infraestrutura e integração com outras áreas do
produto. Essa posição exigiu dar mais atenção à comunicação com o time e maior cuidado
com decisões que poderiam afetar o andamento de vários colegas, além disso como Ages III também
represento uma influência para meus colegas em estágios anteriores do curso, portanto, mais do que nunca
é vital que eu demonstre um bom exemplo, para poder ver colegas se desenvolvendo do mesmo modo.

Do ponto de vista técnico, considero o projeto relevante por envolver
integrações mais complexas, como Firebase, Mapbox, \textit{sockets} e observabilidade,
além de tarefas como \textit{deploy} e organização de ambiente. Também foi uma oportunidade de consolidar
conhecimentos de disciplinas como Construção de Software e Projeto e Arquitetura de Software, principalmente a segunda,
da qual tirei muita inspiração em aulas e principalmente do livro Arquitetura Limpa de Robert Martin, para me ajudarem nas 
descisões arquiteturais implementadas, permitindo balancear um alto nível de desacomplamento ao mesmo tempo que mantendo a 
complexidade de escrever o código relativamente baixa.

Considero positivo ter assumido responsabilidades mais difíceis e participado ativamente da estabilização do projeto
para a entrega. A experiência mostrou que a atuação de um Ages III exige equilíbrio
entre desenvolver, orientar, revisar e remover bloqueios para o time. Não considero que minha atuação foi perfeita, no 
entanto fico com dificuldade em apontar pontos negativos que foram relevantes, creio que o principal foi a demora para 
aproximação com os Ages I e II, que mesmo tarde, em meados da sprint 2, aconteceu e sinceramente a ocorrência tardia não parece 
ter afetado tanto o desempenho do projeto.

Por fim, creio que o VanGO foi um projeto muito importante para minha evolução pessoal e profissional, pois aproximou
minha atuação de desafios mais próximos aos de liderança técnica. Saio desta Ages muito melhor do que entrei, 
com mais maturidade para avaliar riscos, documentar decisões, organizar prioridades
e apoiar colegas em tarefas complexas. Os aprendizados aqui adquiridos servirão como base tanto para
essa última experiência na Ages, quanto para minha formação profissional.
